% =====================================================================
% 第八讲 (Exposé VIII) 忠实平坦下降
% (Descente fidèlement plate)
% M. Raynaud
% 原书页码: pp. 153-176
% =====================================================================
\setcounter{chapter}{7}
\chapter{忠实平坦下降}

\section{拟凝聚模的下降}

设 $\mathbf{Sch}$ 是预概形范畴\index[notations]{$\mathbf{Sch}$}. 如 VI 11 b) 那样处理，我们发现由预概形 $X$ 和 $X$ 上的模 $F$ 组成的对 $(X, F)$ 的范畴 (其态射如 loc. cit. 那样借助模经环化空间态射的直像概念定义) 可以看作 $\mathbf{Sch}$ 之上的纤维范畴，相对 $\mathbf{Sch}$ 中态射 $f \colon X \to Y$ 的基变换函子是模经 $f$ 的逆像函子. (注意前述纤维范畴在 $X \in \mathrm{Ob}(\mathbf{Sch})$ 处的纤维范畴是 $X$ 上模范畴的对偶范畴). 由于拟凝聚模的逆像拟凝聚，看到对 $(X, F)$ 中 $F$ 拟凝聚的配对组成的范畴的全子范畴是前述纤维范畴的纤维子范畴. (相反，若对 $f$ 不作假设，拟凝聚模的直像一般不拟凝聚). 我们简称为预概形上拟凝聚模的纤维范畴.

另一方面回忆，环化空间态射 $f \colon X \to Y$ 称为忠实平坦，若它平坦 (i.e. 对每个 $x \in X$，$\mathcal{O}_{X,x}$ 是 $\mathcal{O}_{Y,f(x)}$ 上的平坦模，(IV))，且满射. 称 $f$ 是\emph{拟紧态射}\index[term]{nijintaishe@拟紧态射}，若任何拟紧部分的原像是拟紧的；当 $f$ 是预概形态射时，这也意味着 $Y$ 的仿射开集的原像是 $X$ 的有限个仿射开集的并.

\begin{theorem}\label{thm:VIII.1.1}
设 $\mathcal{F}$ 是预概形上拟凝聚模的纤维范畴. 设 $g \colon S' \to S$ 是预概形态射，忠实平坦且拟紧. 则 $g$ 是有效 $\mathcal{F}$-下降态射\index[term]{youxiaoxiajiangtaishe@有效下降态射}.
\end{theorem}

回忆\footnote{这里承认下降的一般理论，它在 J. Giraud 的文章 (引自 Avertissement 的脚注 (5)) 中有详细阐述，此后我们把它引用为 \cite{exp08-D}. 也可参见 \cite{exp08-2} 的简明阐述.} 这意味着两件事：

\begin{corollary}\label{cor:VIII.1.2}
(\emph{模同态的下降})\index[term]{motongtaidexiajiang@模同态的下降} 设 $g \colon S' \to S$ 是预概形态射，忠实平坦且拟紧，$F$ 和 $G$ 是 $S$ 上的两个拟凝聚模，$F'$ 和 $G'$ 是它们在 $S'$ 上的逆像，最后 $F''$ 和 $G''$ 是它们在 $S'' = S' \times_S S'$ 上的逆像. 考虑由经 $g$，$p_1$，$p_2$ 的基变换函子 (其中 $p_1, p_2 \colon S' \times_S S' \rightrightarrows S'$ 是两个投影) 定义的集合映射图式：
\[
\mathrm{Hom}_S(F, G) \longrightarrow \mathrm{Hom}_{S'}(F', G') \rightrightarrows \mathrm{Hom}_{S''}(F'', G'').
\]
该图式正合，i.e. 定义第一个集合到第二个到第三个的两个映射的一致集的双射.
\end{corollary}

换言之，经 $g$ 的基变换函子 $F \mapsto F'$ 定义从 $S$ 上拟凝聚模范畴到配备相对 $g$ 的下降数据的 $S'$ 上拟凝聚模范畴的忠实全函子. 此外：

\begin{corollary}\label{cor:VIII.1.3}
(\emph{模的下降})\index[term]{modexiajiang@模的下降} 对 $S'$ 上的每个拟凝聚模 $F'$，$F'$ 上相对 $g$ 的每个下降数据有效，i.e. $F'$ 连同其下降数据同构于 $S$ 上某拟凝聚模经 $g$ 的逆像 (由 \ref{cor:VIII.1.2} 在唯一同构意义下确定).
\end{corollary}

换言之，前述忠实全函子甚至是等价. 实际上，这意味着给出 $S$ 上的拟凝聚模，与给出 $S'$ 上配备相对 $g$ 的下降数据的拟凝聚模，是同一回事.

\begin{proof}
先设 $T$ 是 $S$-预概形，它 $S$-同构于覆盖 $S$ 的诱导开集 $S_i$ 族 $S$ 的和. 则结构态射 $T \to S$ 显然是有效 $\mathcal{F}$-下降态射 (这精确意味着 $S$ 上拟凝聚模 $F$ 的给出等价于诸 $S_i$ 上拟凝聚模 $F_i$ 及满足众所周知的上链条件的粘合同构 $\varphi_{ji} \colon F_i|_{S_i \cap S_j} \to F_j|_{S_i \cap S_j}$ 的给出). 由 VII，8 得出，为验证 $g \colon S' \to S$ 是有效 $\mathcal{F}$-下降态射，只需对由 $g$ 经基变换 $T \to S$ 得到的态射 $g_T \colon T' = T \times_S S' \to T$ 验证. (注意对 $T \to S$ 的假设在任何基变换下保持，故 $T \to S$ 实际上是万有有效 $\mathcal{F}$-下降态射). 取覆盖 $S$ 的仿射开集作 $S_i$，于是归结到 $S$ 仿射的情形.

则 $S'$ 是有限个仿射开集的并，取这些开集的和 $S$-概形，得到仿射 $S$-概形 $S_1$ 和满射平坦的 $S$-态射 $S_1 \to S'$. 故 $S_1$ 在 $S$ 上也忠实平坦. 因此若证明忠实平坦且仿射的态射是有效 $\mathcal{F}$-下降态射，从而是万有严格 $\mathcal{F}$-下降态射 (该假设确实在基变换下保持)，则特别得出结论：结构态射 $S_1 \to S$ 是万有严格 $\mathcal{F}$-下降态射，且由于存在 $S$-态射 $S_1 \to S'$，由 \cite{exp08-D} 将得出 $g \colon S' \to S$ 是严格 $\mathcal{F}$-下降态射.

这归结到 $g$ 是仿射态射的情形，且如所见此时还可假设 $S$ 仿射，故可假设 $S$ 和 $S'$ 仿射. 此时 \ref{cor:VIII.1.2} 等价于
\end{proof}

\begin{lemma}\label{lem:VIII.1.4}
设 $A$ 是环，$A'$ 是忠实平坦的 $A$-代数，$M$ 和 $N$ 是两个 $A$-模，$M'$ 和 $N'$ 是由环变换 $A \to A'$ 得到的 $A'$-模，$M''$，$N''$ 是由环变换 $A \to A''$ 得到的 $A'' = A' \otimes_A A'$-模. 则集合映射序列
\[
\mathrm{Hom}_A(M, N) \longrightarrow \mathrm{Hom}_{A'}(M', N') \rightrightarrows \mathrm{Hom}_{A''}(M'', N'')
\]
正合.
\end{lemma}

由于 $N \to N'$ 单射 ($A'$ 在 $A$ 上忠实平坦)，看到第一个箭单射. 剩下证明：若 $A'$-同态 $u' \colon M' \to N'$ 与下降数据相容，则它来自 $A$-同态 $u \colon M \to N$. 这也简单意味着 $u'$ 把 $M'$ 的子集 $M$ 映入 $N'$ 的子集 $N$ (诱导映射 $u \colon M \to N$ 将自动 $A$-线性，因为 $u'$ 是 $A'$-线性的，同样看到 $u'$ 必然等于 $u \otimes_A A'$). 若 $x \in M$，则 $u'(x)$ 是映射对 $N' \rightrightarrows N''$ 的核中的元素. 为证 \ref{lem:VIII.1.4}，于是归结到下面的特殊情形 (对应于 $M = A$ 的情形)：

\begin{corollary}\label{cor:VIII.1.5}
设 $N$ 是 $A$-模，则集合映射序列
\[
N \longrightarrow N' \rightrightarrows N''
\]
正合.
\end{corollary}

设 $A_1$ 是忠实平坦的 $A$-代数. 为证所述序列正合，只需证经环变换 $A \to A_1$ 得到的序列正合. 后者显然正是 $A_1$-模 $N_1 = N \otimes_A A_1$ 和 $A_1$-代数 $A'_1 = A_1 \otimes_A A'$ 所对应的序列. 故只需找 $A$ 上忠实平坦的 $A_1$，使 $\mathrm{Spec}(A'_1) \to \mathrm{Spec}(A_1)$ 是严格 $\mathcal{F}$-下降态射. 取 $A_1 = A'$ 即可，因为此时前述态射有右逆态射，从而由 \cite{exp08-D} 它对 $\mathbf{Sch}$ 上的任何纤维范畴是有效下降态射.

最后剩下证明：若 $N'$ 是配备对 $A \to A'$ 的下降数据的 $A'$-模，i.e. 配备由 $N$ 经环变换 $A' \rightrightarrows A' \otimes_A A'$ 得到的两个模之间的同构
\[
\varphi \colon N'_1 \xrightarrow{\sim} N'_2,
\]
则 $N'$ 连同其下降数据同构于模 $N \otimes_A A'$. 考虑到 \ref{cor:VIII.1.5}，容易看到该断言等价于：

\begin{lemma}\label{lem:VIII.1.6}
设 $N'$ 是配备相对 $A \to A'$ (其中 $A'$ 是 $A$-代数) 的下降数据的 $A'$-模. 设 $N$ 是 $N'$ 的由满足 $\varphi(x \otimes_{A'} 1_{A'}) = 1_{A'} \otimes_{A'} x$ 的 $x$ 组成的子 $A$-模，考虑典范同态
\[
N \otimes_A A' \longrightarrow N',
\]
(它与下降数据相容). 若 $A'$ 在 $A$ 上忠实平坦，则该同态是同构.
\end{lemma}

\begin{proof}
证明该引理. 再设 $A_1$ 是忠实平坦的 $A$-代数，为证所述态射是同构，只需证它在环变换 $A \to A_1$ 后成为同构. 利用 $A_1$ 在 $A$ 上的平坦性，看到如此得到的同态正是按 $A'_1 = A' \otimes_A A_1$ 上的模 $N' \otimes_A A_1$ 项下、配备相对 $A_1 \to A'_1$ 的由 $N'$ 上给定数据经环变换典范得到的下降数据而直接得到的那一个. 故只需找 $A$ 上忠实平坦的 $A_1$，使 $\mathrm{Spec}(A'_1) \to \mathrm{Spec}(A_1)$ 是有效 $\mathcal{F}$-下降态射. 如上取 $A_1 = A'$. 这完成 \ref{lem:VIII.1.6} 的证明，从而完成 \ref{thm:VIII.1.1} 的证明.
\end{proof}

\begin{corollary}\label{cor:VIII.1.7}
(模的截面的下降)\index[term]{mojiemiandexiajiang@模截面的下降} 设 $g \colon S' \to S$ 是预概形态射，忠实平坦且拟紧. 对 $S$ 上的每个拟凝聚模 $G$，设 $G'$ 和 $G''$ 是它在 $S'$ 和 $S'' = S' \times_S S'$ 上的逆像，考虑 $S$ 上模同态图式
\[
G \longrightarrow g_{*}(G') \rightrightarrows h_{*}(G'')
\]
(其中 $h \colon S'' \to S$ 是结构态射). 该图式正合.
\end{corollary}

事实上，这意味着对 $S$ 中的每个开集 $U$，由 $U$ 上截面组成的对应图式正合. 显然可假设 $U = S$，此时所述正合性是 \ref{cor:VIII.1.2} 的特例，取 $F = \mathcal{O}_S$ 得到.

由于模的逆像函子右正合，由 \ref{thm:VIII.1.1} 形式地得出：

\begin{corollary}\label{cor:VIII.1.8}
(\emph{商模的下降})\index[term]{shangmodexiajiang@商模的下降} 用 \ref{cor:VIII.1.7} 的记号，设 $\mathrm{Quot}(F)$ 对每个预概形上的拟凝聚模 $F$ 表示 $F$ 的拟凝聚商模集. 在此约定下，集合映射图式
\[
\mathrm{Quot}(G) \longrightarrow \mathrm{Quot}(G') \rightrightarrows \mathrm{Quot}(G'')
\]
正合.
\end{corollary}

(用子模代替商模显然有同样的断言，因为两者一一对应). 特别取 $G = \mathcal{O}_S$，得到：

\begin{corollary}\label{cor:VIII.1.9}
(\emph{闭子预概形的下降})\index[term]{biziyugaixingdexiajiang@闭子预概形的下降} 对每个预概形 $X$，设 $H(X)$ 表示 $X$ 的闭子预概形集. 在此记号下，在 \ref{cor:VIII.1.7} 的条件下，下列集合映射图式
\[
H(S) \longrightarrow H(S') \rightrightarrows H(S'')
\]
正合.
\end{corollary}

有必要用下列结果补充\cref{thm:VIII.1.1}：

\begin{proposition}\label{prop:VIII.1.10}
(模的性质的下降)\index[term]{moxingzhidexiajiang@模性质的下降} 设 $g \colon S' \to S$ 是忠实平坦拟紧态射，$F$ 是 $S$ 上的拟凝聚模. 为使 $F$ 有限型 (resp. 有限表现，resp. 局部自由且有限型)，必要且充分条件是它在 $S'$ 上的逆像 $F'$ 如此.
\end{proposition}

只需证明"充分"部分. 显然可假设 $S$ 仿射，再把 $S'$ 换成覆盖 $S'$ 的仿射开集的和，归结到 $S'$ 也仿射的情形. 此时断言等价于：

\begin{corollary}\label{cor:VIII.1.11}
设 $A$ 是环，$A'$ 是忠实平坦的 $A$-代数，$M$ 是 $A$-模，$M'$ 是 $A'$-模 $M \otimes_A A'$. 为使 $M$ 有限型 (resp. 有限表现，resp. 有限型局部自由)，必要且充分条件是 $M'$ 如此.
\end{corollary}

事实上，有 $M = \varinjlim_i M_i$，其中 $M_i$ 是 $M$ 的有限型子模. 从而 $M' = \varinjlim_i M'_i$，若 $M'$ 有限型，则 $M'$ 等于某个 $M'_i$，由忠实平坦性 $M$ 等于 $M_i$，故 $M$ 有限型. 于是存在正合序列
\[
0 \longrightarrow R \longrightarrow L \longrightarrow M \longrightarrow 0,
\]
其中 $L$ 有限型自由，从而有正合序列
\[
0 \longrightarrow R' \longrightarrow L' \longrightarrow M' \longrightarrow 0,
\]
其中 $L'$ 有限型自由. 若 $M'$ 有限表现，则 $R'$ 有限型，由前面所述 $R$ 有限型，故 $M$ 有限表现. 最后，说 $M$ 局部自由且有限型，意味着它有限表现且平坦 (参见 IV，诺特情形；一般情形留给读者). 由于这些性质各自都随下降保持，其合取亦然，这完成证明.

\begin{remark}\label{rem:VIII.1.12}
\ref{thm:VIII.1.1} 与 \ref{prop:VIII.1.10} 的合取表明，把纤维范畴 $\mathcal{F}$ 换成由有限型 (resp. 有限表现，resp. 有限型局部自由，resp. 秩为给定 $n$ 的局部自由) 拟凝聚模组成的纤维子范畴时，断言 \ref{thm:VIII.1.1} 仍然成立.
\end{remark}

\section{仿射预概形在另一个之上的下降}

由于模的逆像函子与张量积及其他张量运算相容，\cref{thm:VIII.1.1} 蕴含各种变体，只需考虑(代替单个拟凝聚模)配备各种借助张量运算表达的附加结构的拟凝聚模或拟凝聚模组. 例如，$S$ 上三个拟凝聚模 $F$，$G$，$H$ 及配对
\[
F \otimes G \longrightarrow H
\]
的给出，等价于 $S'$ 上三个拟凝聚模 $F'$，$G'$，$H'$，配备相对 $g \colon S' \to S$ 的下降数据，并配备与这些下降数据"相容" (按此词的显然含义) 的配对
\[
F' \otimes G' \longrightarrow H'
\]
的给出. 例如，若 $F = G = H$，看到 $S$ 上配备代数律 (暂时不假设它满足结合律、交换律或单位截面的存在性公理) 的拟凝聚模 $F$ 的给出，等价于 $S'$ 上同样的给出，外加一个下降数据. 利用前一节的结果，立即看到为使 $F$ 满足刚才提到的通常公理之一，必要且充分条件是 $F'$ 如此. 例如，$S$ 上拟凝聚代数 (此后我们用此词总指：结合的、交换的、有单位截面的) $\mathcal{A}$ 的给出，等价于 $S'$ 上拟凝聚代数 $\mathcal{A}'$ 配备相对 $g \colon S' \to S$ 的下降数据的给出. 若回忆 $S$ 上拟凝聚代数的对偶范畴与 $S$ 上仿射 $S$-预概形范畴之间的等价 (EGA II §1)，立即得到：

\begin{theorem}\label{thm:VIII.2.1}
设 $\mathcal{F}'$ 是预概形仿射态射 $f \colon X \to S$ 的纤维范畴，把它看作预概形范畴 $\mathbf{Sch}$ 中箭的纤维范畴 (VI 11.a) 的纤维子范畴. 设 $g \colon S' \to S$ 是忠实平坦拟紧的预概形态射. 则 $g$ 是有效 $\mathcal{F}'$-下降态射.
\end{theorem}

\section{态射的集合性质与有限性性质的下降}

\begin{proposition}\label{prop:VIII.3.1}
设 $f \colon X \to Y$ 是 $S$-态射\footnote{关于第 3 和第 4 号处理的这类结果的其他结果，参见 EGA IV 2.3，2.6，2.7.}，$g \colon S' \to S$ 是满射态射，$f' \colon X' = X \times_S S' \to Y' = Y \times_S S'$ 是由 $f$ 借助 $g \colon S' \to S$ 作基变换得到的态射. 为使 $f$ 满射 (resp. 根式)，必要且充分条件是 $f'$ 如此.
\end{proposition}

注意 $f'$ 也可经基变换 $Y' \to Y$ 得到，后者也满射，因为它由满射的 $g \colon S' \to S$ 得到. 另一方面，对每个 $y \in Y$ 和 $y$ 上的每个 $y' \in Y'$，有同构
\[
X'_{y'} \xleftarrow{\sim} X_y \otimes_{k(y)} k(y'),
\]
其中 $X_y$ 表示 $X$ 在 $y$ 处的纤维，$X'_{y'}$ 表示 $X'$ 在 $y'$ 处的纤维. 由此得出，$X_y$ 非空 (resp. 至多有一个点，且该点对应根式剩余扩张) 当且仅当 $X'_{y'}$ 有同样性质. 这证明 3.1.

\begin{corollary}\label{cor:VIII.3.2}
在 \ref{prop:VIII.3.1} 的条件下，若 $f'$ 单射 (resp. 双射)，则 $f$ 也如此.
\end{corollary}

这是因为若 $X'_{y'}$ 至多有一个点 (resp. 恰好一个点)，则 $X_y$ 亦然；确实如此，因为态射 $X'_{y'} \to X_y$ 满射 (它由满射的 $\mathrm{Spec}(k(y')) \to \mathrm{Spec}(k(y))$ 得到).

\begin{proposition}\label{prop:VIII.3.3}
用 \ref{prop:VIII.3.1} 的记号. 假设 $g \colon S' \to S$ 满射且拟紧 (resp. 忠实平坦且拟紧). 为使 $f$ 拟紧 (resp. 有限型)，必要且充分条件是 $f'$ 如此.
\end{proposition}

只需证明"充分"部分. 显然可假设 $S = Y$，因为对 $g \colon S' \to S$ 的假设对 $Y' \to Y$ 保持. 还可假设 $Y$ 仿射. 则 $Y'$ 拟紧，故 $X'$ 拟紧 (因为 $f'$ 由假设拟紧). 设 $(X_i)_{i \in I}$ 是覆盖 $X$ 的仿射开集族，则诸 $X'_i$ 是覆盖 $X'$ 的开集，故有覆盖 $X'$ 的有限子族. 由于 $X' \to X$ 满射，得出对应的诸 $X_i$ 已经覆盖 $X$，故 $X$ 拟紧，i.e. $f$ 拟紧. 现在假设 $f'$ 有限型，在 $g$ 忠实平坦的假设下证明 $f$ 有限型. 把 $Y'$ 换成覆盖它的仿射开集族的和，可假设 $Y'$ 仿射. 最后，由前面所述 $X$ 被有限个仿射开集 $X_i$ 覆盖，在已知 $X'_i$ 在 $Y'$ 上有限型的条件下需证明它们在 $Y$ 上有限型. 这归结到：

\begin{corollary}\label{cor:VIII.3.4}
设 $B$ 是 $A$-代数，$A'$ 是忠实平坦的 $A$-代数，$B' = B \otimes_A A'$ 是由 $B$ 经环变换得到的 $A'$-代数. 为使 $B$ 有限型，必要且充分条件是 $B'$ 如此.
\end{corollary}

只需证明"充分"部分. 有 $B = \varinjlim_i B_i$，其中 $B_i$ 是 $B$ 的有限型子代数. 故有 $B' = \varinjlim_i B'_i$，若 $B'$ 在 $A'$ 上有限型，则 $B'$ 等于某个 $B'_i$，故 $B$ 等于 $B_i$，从而有限型.

\begin{corollary}\label{cor:VIII.3.5}
仍假设基变换态射 $g \colon S' \to S$ 忠实平坦且拟紧. 为使 $f$ 拟有限，必要且充分条件是 $f'$ 如此.
\end{corollary}

事实上，"拟有限"性质按定义是"有限型"与"纤维有限"的合取，二者各自都能随 $g$ 下降，第一个由 \ref{prop:VIII.3.3}，第二个由 \ref{prop:VIII.3.1} 的推理 (只用 $g$ 满射的事实).

\begin{remark}\label{rem:VIII.3.6}
设 $A$ 是环，$X$ 是 $A$-预概形. 容易看到下列条件等价：
\begin{enumerate}[i)]
\item 存在诺特环 $A_0$ (若愿意可假设它是 $A$ 的有限型子环)，有限型的 $A_0$-预概形 $X_0$，同态 $A_0 \to A$ 和 $A$-同构 $X \xrightarrow{\sim} X_0 \times_{A_0} A$.
\item 对角态射 $X \to X \times_{\mathrm{Spec}(A)} X$ 拟紧 (若 $X$ 在 $A$ 上分离则此条件为空的)，$X$ 是有限个仿射开集 $X_i$ 的并，其环 $B_i$ 是 $A$ 上有限表现的代数，i.e. 有限个未定元的多项式代数被有限型理想的商.
\end{enumerate}
若 $X$ 本身仿射，环为 $B$，这些条件简单意味着 $B$ 是 $A$ 上有限表现的代数.

态射 $f \colon X \to Y$ 称为\emph{有限表现态射}\index[term]{youxianbiaoxiantaishe@有限表现态射}，也说 $X$ 在 $Y$ 上有限表现，若 $Y$ 是仿射开集 $Y_i$ 的并，使 $X|_{Y_i}$ 作为 $Y_i$-预概形满足前面等价的条件. 则对 $Y$ 中每个仿射开集 $Y'$，$X|_{Y'}$ 亦然. 这是随基变换保持的性质，而且两个有限表现态射的复合有限表现.

这些概念确定后，由 (ii) 看到 (如 \ref{prop:VIII.1.10} 那样进行) 把其中的"有限型"换成"有限表现"，该断言仍然成立.
\end{remark}

\section{拓扑性质的下降}

\begin{theorem}\label{thm:VIII.4.1}
设 $g \colon Y' \to Y$ 是态射，$Z$ 是 $Y$ 的部分. 假设 $g$ 平坦，且存在拟紧态射 $f \colon X \to Y$ 使 $Z = f(X)$ (N.B. 若 $Y$ 诺特，这最后一个条件由"$Z$ 可构"蕴含). 则
\[
\overline{g^{-1}(Z)} = g^{-1}(\overline{Z}).
\]
\end{theorem}

可假设 $Y$ 仿射，再假设 $Y'$ 仿射. 由于 $Y$ 仿射，$X$ 是有限个仿射开集 $X_i$ 的并，把 $X$ 换成诸 $X_i$ 的和，可假设 $X$ 也仿射. 设 $A$，$A'$，$B$ 是 $Y$，$Y'$，$X$ 的环，$B' = B \otimes_A A'$ 是 $X' = X \times_Y Y'$ 的环，$I$ 是 $A \to B$ 的核，$I'$ 是 $A' \to B'$ 的核，故由这些理想定义的 $Y$ 和 $Y'$ 的闭部分分别是 $Z = f(X)$ 的闭包和 $Z' = f'(X') = g^{-1}(Z)$ 的闭包. 要证明后者等于 $g^{-1}(\overline{Z})$，这将由 $I' = IA'$ 得出，而后者又是 $A'$ 在 $A$ 上平坦的推论.

\begin{corollary}\label{cor:VIII.4.2}
设 $g \colon Y' \to Y$ 是平坦拟紧态射，$Z'$ 是 $Y'$ 的对于由 $g$ 定义的集合等价关系饱和的闭部分. 则 $Z' = g^{-1}(g(Z'))$.
\end{corollary}

确实有 $Z' = g^{-1}(Z)$，其中 $Z = g(Z')$. 可以应用 \ref{thm:VIII.4.1}，注意 \ref{thm:VIII.4.1} 中对 $Z$ 的条件确实满足，只需取 $X$ 为配备由 $Y'$ 诱导的既约结构的预概形 $Z'$. ($g$ 拟紧保证诱导态射 $f \colon Z' \to Y$ 拟紧).

断言 \ref{cor:VIII.4.2} 也意味着由 $Y$ 诱导的 $g(Y')$ 的拓扑是 $Y'$ 的拓扑的商. 特别：

\begin{corollary}\label{cor:VIII.4.3}
设 $g \colon Y' \to Y$ 是忠实平坦拟紧态射. 则 $g$ 使 $Y$ 成为 $Y'$ 的商拓扑空间，i.e. 对 $Y$ 的部分 $Z$，$Z$ 闭 (resp. 开) 当且仅当 $Z' = g^{-1}(Z)$ 如此.
\end{corollary}

回忆 $Y'$ 的两个元素 $a$，$b$ 在 $Y$ 中有相同的像，当且仅当它们形如 $p_1(c)$，$p_2(c)$，其中 $c$ 是 $Y'' = Y' \times_Y Y'$ 中适当的元素. 由此得出，若 $g$ 满射，有集合的精确图式
\[
P(Y) \longrightarrow P(Y') \rightrightarrows P(Y''),
\]
其中对每个集合 $E$，用 $P(E)$ 表示它的部分集. 这样，\ref{cor:VIII.4.3} 也可解释为：

\begin{corollary}\label{cor:VIII.4.4}
(开部分的下降 resp. 闭部分的下降) 设 $g \colon Y' \to Y$ 如 \ref{cor:VIII.4.3} 中. 对每个预概形 $X$，设 $\mathrm{Ouv}(X)$ (resp. $\mathrm{Fer}(X)$) 表示它的开部分集 (resp. 闭部分集). 则有集合映射的精确图式 (由 $g$ 和 $Y'' = Y' \times_Y Y'$ 的两个投影得到)：
\[
\mathrm{Ouv}(Y) \longrightarrow \mathrm{Ouv}(Y') \rightrightarrows \mathrm{Ouv}(Y''),
\]
\[
\mathrm{Fer}(Y) \longrightarrow \mathrm{Fer}(Y') \rightrightarrows \mathrm{Fer}(Y'').
\]
\end{corollary}

有下列对 \ref{cor:VIII.4.3} 的补充：

\begin{corollary}\label{cor:VIII.4.5}
设 $g \colon Y' \to Y$ 如 \ref{cor:VIII.4.3} 中，设 $Z$ 是 $Y$ 的以某个拟紧态射 $f \colon X \to Y$ 的像为 $Z$ 的部分 (例如 $Z$ 可构，$Y$ 诺特). 为使 $Z$ 是 $Y$ 的局部闭部分，必要且充分条件是 $Z' = g^{-1}(Z)$ 是 $Y'$ 的局部闭部分.
\end{corollary}

只需证明"充分"部分. 设 $Y_1$ 是 $Y$ 的闭子预概形，即 $Z$ 的闭包配备诱导既约结构，设 $Y'_1 = Y_1 \times_Y Y'$ 是 $Y'$ 的 $Y_1$ 的原像构成的闭子预概形. 它的底集是原像 $g^{-1}(Y_1) = g^{-1}(\overline{Z})$，由 \ref{thm:VIII.4.1} 等于 $Z'$. 由于 $Z'$ 在 $Y'$ 中局部闭，它在 $Z'$ 中开，故在 $Y'_1$ 中开. 而后者在 $Y_1$ 上忠实平坦且拟紧，由 \ref{cor:VIII.4.3} 得出 $Z$ 在 $Y_1$ 中开，i.e. 在 $\overline{Z}$ 中开，这意味着 $Z$ 局部闭.

\begin{corollary}\label{cor:VIII.4.6}
设 $g \colon S' \to S$ 是忠实平坦拟紧态射，$f \colon X \to Y$ 是 $S$-态射，$f' \colon X' \to Y'$ 是经基变换得到的 $S'$-态射. 假设 $f'$ 是开映射 (resp. 闭映射，resp. 拟紧且嵌入同胚，resp. 满同胚)；则 $f$ 有同样性质.
\end{corollary}

由于 $Y'$ 在 $Y$ 上忠实平坦且拟紧，可假设 $Y = S$. 设 $Q$ 是 $X$ 的部分，有 (用 $h$ 记投影态射 $X' \to X$)：
\[
g^{-1}(f(Q)) = f'(h^{-1}(Q)).
\]
若 $Q$ 开 (resp. 闭)，则 $h^{-1}(Q)$ 亦然，故假设 $f'$ 是开映射 (resp. 闭映射) 时 $f'(h^{-1}(Q))$ 亦然，由前面公式和 \ref{cor:VIII.4.3} 得 $f(Q)$ 亦然. 这证明 \ref{cor:VIII.4.6} 中前两个断言，剩下考察 $f'$ 是嵌入同胚的情形，证明 $f$ 是嵌入同胚. (满同胚的情形将由 \ref{prop:VIII.3.1} 得出). 由 \ref{cor:VIII.3.2} $f$ 单射，剩下证明映射 $X \to f(X)$ 开. 已知由 \ref{prop:VIII.3.3} $f$ 拟紧. 只需证明对 $X$ 的每个闭部分 $Z$，有 $Z = f^{-1}(\overline{f(Z)})$，这等价于满射映射 $h \colon X' \to X$ 的逆像的类似公式，i.e. 等价于
\[
Z' = f'^{-1}(g^{-1}(\overline{f(Z)})),
\]
其中令 $Z' = h^{-1}(Z)$. 而由 \ref{thm:VIII.4.1} 应用于 $Y$ 的部分 $f(Z)$，有 $g^{-1}(\overline{f(Z)}) = \overline{g^{-1}(f(Z))}$，待证公式等价于
\[
Z' = f'^{-1}(\overline{f'(Z')}),
\]
它由 $f'$ 是嵌入同胚的假设得出.

N.B. 在后一个推理中，已经假设 $f$ 拟紧，没有用到 $g$ 拟紧，只用 $g$ 忠实平坦. 故正是在这个假设下可以下降"嵌入同胚"或"满同胚"性质，或借助前面推理下降"$f'$ 拟紧且使 $f'(X')$ 成为 $X'$ 的商拓扑空间"的性质.

我们说预概形态射 $f \colon X \to Y$ 万有开 (resp. 万有闭，resp. 万有双连续，等等)，若对任何基变换 $Y' \to Y$，$f' \colon X' \to Y'$ 是开态射 (resp. 闭态射，resp. 到像空间的同胚). 由 \ref{cor:VIII.4.6} 得出：

\begin{corollary}\label{cor:VIII.4.7}
在 \ref{cor:VIII.4.6} 的条件下，为使 $f$ 万有开 (resp. 万有闭，resp. 万有嵌入同胚，resp. 万有同胚)，必要且充分条件是 $f'$ 如此.
\end{corollary}

\begin{corollary}\label{cor:VIII.4.8}
在 \ref{cor:VIII.4.6} 的条件下，为使 $f$ 分离 (resp. 紧合)，必要且充分条件是 $f'$ 如此.
\end{corollary}

说 $f$ 分离意味着对角态射 $X \to X \times_Y X$ 闭，或也万有闭，故 \ref{cor:VIII.4.8} 的第一个断言由 \ref{cor:VIII.4.7} 得出. 说 $f$ 紧合意味着 $f$ 满足条件 a) $f$ 有限型 b) $f$ 分离 c) $f$ 万有闭. 条件 a) 由 \ref{prop:VIII.3.3} 下降良好，b) 由刚才所述亦然，最后 c) 由 \ref{cor:VIII.4.7} 亦然.

\begin{remark}\label{rem:VIII.4.9}
回忆当 $g \colon Y' \to Y$ 是有限型平坦态射，$Y$ 局部诺特时，$g$ 是开态射 (\cref{prop:VI.6.6})，这是比 \ref{cor:VIII.4.3} 更精确的结果. 然而注意，若 $f$ 是诺特预概形的忠实平坦拟紧态射，$f$ 一般不一定是开态射. 例如设 $Y$ 是生成点 $y$ 不开的不可约概形 (例如代数曲线)，取 $Y'$ 为概形和 $Y \amalg \mathrm{Spec}(k(y))$，则开部分 $\mathrm{Spec}(k(y))$ 经结构态射 $Y' \to Y$ 的像不是 $Y$ 的开部分. 读者还会注意，本讲的各种断言若放弃所考虑的忠实平坦态射还拟紧的假设将变假，使断言失效的典型情形是把 $Y'$ 取为 $Y$ 诸点的局部环的谱的概形和. 例如，再取 $Y$ 为不可约代数曲线，$Z$ 为约化为生成点的部分，它在 $Y'$ 中的逆像开，而 $Z$ 不开.
\end{remark}

4.10. 本讲的各种断言中，把"$Y'$ 在 $Y$ 上平坦"的假设换成：存在 $Y'$ 上有限型、支集为 $Y'$、相对 $Y$ 平坦的模 $F$；忠实平坦假设则换成前述假设加上 $Y' \to Y$ 满射. 这适用于 \ref{prop:VIII.1.10} 中前两个断言，\ref{prop:VIII.3.3}，\ref{cor:VIII.3.5}，\ref{thm:VIII.4.1}，从而适用于本节所有结果.

\section{预概形态射的下降}

\begin{proposition}\label{prop:VIII.5.1}
设 $g \colon S' \to S$ 是预概形态射.
\begin{enumerate}[a)]
\item 假设 $g$ 满射，且同态
\[
g^{*} \colon \mathcal{O}_S \longrightarrow g_*(\mathcal{O}_{S'})
\]
单射，则 $g$ 是预概形范畴中的满态射，甚至在环化空间范畴中亦然.
\item 假设 $g$ 满射且使 $S$ 成为 $S'$ 的商拓扑空间. 设 $S'' = S' \times_S S'$，设 $h \colon S'' \to S$ 是结构态射，考虑典范同态图式
\[
\mathcal{O}_S \longrightarrow g_*(\mathcal{O}_{S'}) \rightrightarrows h_*(\mathcal{O}_{S''}),
\]
并假设该图式正合. 则 $g$ 是预概形范畴中的\emph{有效满态射} (也在环化空间范畴中)，i.e. 图式
\[
S \longleftarrow S' \leftleftarrows S''
\]
正合.
\end{enumerate}
\end{proposition}

\begin{proof}
a) 需证明环化空间态射 $f \colon S \to Z$ 在已知 $fg$ 时确定. 由于 $g$ 满射，已知 $f$ 的底集映射 $f_0$，剩下确定环层同态 $\mathcal{O}_Z \to \mathcal{O}_S$，或等价地由 $f$ 定义的环层同态
\[
u \colon f_0^{-1}(\mathcal{O}_Z) \longrightarrow \mathcal{O}_S.
\]
已知由 $fg$ 定义的同态
\[
(fg)_0^{-1}(\mathcal{O}_Z) = g_0^{-1}(f_0^{-1}(\mathcal{O}_Z)) \longrightarrow \mathcal{O}_{S'},
\]
或等价地，有一个同态
\[
f_0^{-1}(\mathcal{O}_Z) \longrightarrow g_{0*}(\mathcal{O}_{S'}) = g_*(\mathcal{O}_{S'}).
\]
立即看到后者正是 $g^{*} \colon \mathcal{O}_S \to g_*(\mathcal{O}_{S'})$ 与 $u$ 的复合，且由于 $g^{*}$ 单射，已知 $g^{*} \circ u$ 时 $u$ 确定. [N.B. 显然没有用到 $g \colon S' \to S$ 是预概形态射，断言对任何环化空间态射成立；同样的评注对 b) 成立，无论在环化空间范畴中，还是在以局部环为环的环化空间范畴中. 还注意若 $g$ 是预概形态射，未必满射，但 $g^{*} \colon \mathcal{O}_S \to g_*(\mathcal{O}_{S'})$ 单射，则对 $S$ 到概形 $Z$ 的两个态射 $f_1$，$f_2$，由 $f_1 g = f_2 g$ 得 $f_1 = f_2$；事实上，若 $I$ 是 $S$ 上定义 $f_1$，$f_2$ 的一致点子预概形的理想 (即 $Z \times Z$ 的对角子预概形经 $(f_1, f_2)$ 的原像)，看到 $I$ 包含在 $\mathrm{Ker}(g^{*})$ 中].

b) 需证明对每个环化空间 $Z$，下列映射图式
\[
\mathrm{Hom}(S, Z) \longrightarrow \mathrm{Hom}(S', Z) \rightrightarrows \mathrm{Hom}(S'', Z)
\]
正合，且当 $Z$ 是以局部环为环的环化空间、限于以局部环为环的环化空间同态时亦然. 由 a) 已知第一个映射单射，剩下看：若 $f' \colon S' \to Z$ 是环化空间同态且 $f' p_1 = f' p_2$，则 $f'$ 形如 $f g$，其中 $f \colon S \to Z$ 是环化空间同态. 由于 $g$ 满射，显然若 $f'$ 是局部环的环化空间态射，则 $f$ 亦然.

由对 $f'$ 的假设得出底集映射 $f'_0$ 在映射 $g_0$ 的诸纤维上为常值，故由于后者满射，$f'_0$ 唯一分解为 $f'_0 = f_0 g_0$，其中 $f_0 \colon S \to Z$ 是映射，它必然连续，因为 $g_0$ 把 $S$ 等同于 $S'$ 的商拓扑空间. 现在考虑由对应 $f'$ 的同态 $(f_0 g_0)^{-1}(\mathcal{O}_Z) \to \mathcal{O}_{S'}$ 得到的同态
\[
f_0^{-1}(\mathcal{O}_Z) \longrightarrow g_*(\mathcal{O}_{S'}).
\]
假设 $f' p_1 = f' p_2$ 解释为：前述同态与两个同态 $g_*(\mathcal{O}_{S'}) \rightrightarrows h_*(\mathcal{O}_{S''})$ 的复合相同，故由假设 b) 它经一个态射
\[
f_0^{-1}(\mathcal{O}_Z) \longrightarrow \mathcal{O}_S
\]
分解. 后者定义环化空间态射 $f \colon S \to Z$，即所求态射.
\end{proof}

\begin{theorem}\label{thm:VIII.5.2}
设 $\mathcal{F}$ 是预概形范畴 $\mathbf{Sch}$ 中箭的纤维范畴 (VI 11.a). 则任何忠实平坦拟紧态射 $g \colon S' \to S$ 是 $\mathcal{F}$-下降态射 (或如常所说的，$\mathbf{Sch}$ 中的下降态射).
\end{theorem}

这即意味着：设 $S'' = S' \times_S S'$，对 $S$ 上的两个预概形 $X$，$Y$，考虑它们在 $S'$ 上的逆像 $X'$，$Y'$ 及在 $S''$ 上的逆像 $X''$，$Y''$，得映射图式
\[
\mathrm{Hom}_S(X, Y) \longrightarrow \mathrm{Hom}_{S'}(X', Y') \rightrightarrows \mathrm{Hom}_{S''}(X'', Y'');
\]
记号确定后，\ref{thm:VIII.5.2} 意味着该图式正合. 注意一般并不成立 $g$ 是有效下降态射，i.e. 对 $S'$ 上的每个预概形 $X'$，$X'$ 上相对 $g \colon S' \to S$ 的每个下降数据有效. 有效性的问题 (常常微妙) 将在第 7 号考察.

已见 \cite{exp08-D} (考虑到 $\mathbf{Sch}$ 中纤维积存在) 断言 \ref{thm:VIII.5.2} 等价于：

\begin{corollary}\label{cor:VIII.5.3}
预概形的忠实平坦拟紧态射是万有有效满态射.
\end{corollary}

由于忠实平坦拟紧态射经任何基扩张后仍如此，归结到证明它是有效满态射\index[term]{youxiaomantaishe@有效满态射}. 应用判则 \ref{prop:VIII.5.1} b)，它考虑到 \ref{cor:VIII.4.3} 和 \ref{cor:VIII.1.7} 给出所求结果.

\begin{corollary}\label{cor:VIII.5.4}
设 $g \colon S' \to S$ 是忠实平坦拟紧态射，$f \colon X \to Y$ 是 $S$-态射，$f' \colon X' \to Y'$ 是经基变换 $S' \to S$ 得到的 $S'$-态射. 为使 $f$ 是同构，必要且充分条件是 $f'$ 如此.
\end{corollary}

事实上，若 $f'$ 是同构，它也是 $X'$，$Y'$ 上自然下降结构的同构，且由于函子 $X \mapsto X'$ (从 $\mathbf{Sch}_{/S}$ 到配备相对 $g$ 的下降数据的 $S'$ 上预概形范畴) 由 \ref{thm:VIII.5.2} 忠实全，得到所求结论.

\begin{corollary}\label{cor:VIII.5.5}
在 \ref{cor:VIII.5.4} 的条件下，为使 $f$ 是闭浸入 (resp. 开浸入，resp. 拟紧浸入)，必要且充分条件是 $f'$ 如此.
\end{corollary}

如常可假设 $Y = S$，只需证明"充分"部分. 注意 $X'/Y'$ 配备相对 $g \colon Y' \to Y$ 的下降数据，且结构态射 $f' \colon X' \to Y'$ 是浸入故是单态射，这蕴含 $Y''$ 的由 $S''$ 到 $S'$ 的两个投影之一另一个对 $X'/Y'$ 取逆像得到的两个子对象相同. 若 $f'$ 是闭浸入，由 \ref{cor:VIII.1.9} 得出存在 $Y$ 的闭子预概形 $X_1$，其经 $g \colon Y' \to Y$ 的逆像是 $X'$. 故由 $\mathcal{F}$-下降态射相对下降问题的解的唯一性，得出 $X_1$ 与 $X$ $Y$-同构，故 $f \colon X \to Y$ 是闭浸入. 对开浸入用 \ref{cor:VIII.4.4} 同样进行. 最后若 $f'$ 是拟紧浸入，由 \ref{prop:VIII.3.3} $f$ 拟紧，可以对 $Y$ 的部分 $f(X)$ 应用判则 \ref{cor:VIII.4.5}，它证明 $f(X)$ 局部闭，因为它在 $Y'$ 中的逆像 $f'(X')$ 如此. 再把 $Y$ 换成 $f(X)$ 在其中闭的开部分，归结到 $f'$ 是闭浸入的情形，故由前面所述 $f$ 如此.

\begin{corollary}\label{cor:VIII.5.6}
在 \ref{cor:VIII.5.4} 的条件下，为使 $f$ 仿射，必要且充分条件是 $f'$ 如此.
\end{corollary}

如 \ref{cor:VIII.5.5} 那样进行，用 \ref{thm:VIII.2.1} (也可用 Serre 的上同调判则 [EGA II 5.2]，它不借助下降技巧证明 \ref{cor:VIII.5.6}).

\begin{corollary}\label{cor:VIII.5.7}
在 \ref{cor:VIII.5.4} 的条件下，为使 $f$ 整 (resp. 有限，resp. 有限且局部自由)，必要且充分条件是 $f'$ 如此.
\end{corollary}

只需证明"充分"部分，如常可假设 $Y = S$，$Y$ 仿射，$Y'$ 仿射. 由于假设蕴含 $f'$ 仿射，由 \ref{cor:VIII.5.6} $f$ 也仿射，故 $X$ 从而 $X'$ 仿射. 设 $A$，$A'$，$B$，$B' = B \otimes_A A'$ 是 $Y$，$Y'$，$X$，$X'$ 的环. 有 $B = \varinjlim_i B_i$，其中 $B_i$ 遍历 $B$ 的在 $A$ 上有限型的子 $A$-代数，从而 $B' = \varinjlim_i B'_i$，其中 $B'_i$ 是 $A'$-代数 $B'$ 的有限型子代数. 若 $B'$ 在 $A'$ 上整，诸 $B'_i$ 是 $A'$ 上有限型模，故 $A'$ 在 $A$ 上忠实平坦时，诸 $B_i$ 是 $A$ 上有限型模，i.e. $B$ 在 $A$ 上整. 同样看到若 $B'$ 在 $A'$ 上有限，则 $B$ 在 $A$ 上有限. 对"有限型局部自由"同样结论，参见 1.11.

\begin{corollary}\label{cor:VIII.5.8}
在 \ref{cor:VIII.5.4} 的条件下，假设 $f$ 拟紧，设 $\mathcal{L}$ 是 $X$ 上的可逆模，$\mathcal{L}'$ 是它在 $X'$ 上的逆像. 为使 $\mathcal{L}$ 相对 $f$ 丰富 (resp. 非常丰富)，必要且充分条件是 $\mathcal{L}'$ 相对 $f'$ 丰富 (resp. 非常丰富).
\end{corollary}

只需证明"充分"部分. 对 $\mathcal{L}$ 的假设无论如何蕴含 $f'$ 分离，故由 \ref{cor:VIII.4.8} $f$ 分离，且由于 $f$ 拟紧，$g \colon Y' \to Y$ 平坦，经仿射覆盖计算直像表明有同构
\[
g^{*}(f_*(\mathcal{L}^{\otimes n})) \xrightarrow{\sim} f'_*(\mathcal{L}'^{\otimes n})
\]
对每个整数 $n$，故有同构
\[
g^{*}(\mathcal{S}) \xrightarrow{\sim} \mathcal{S}',
\]
其中 $\mathcal{S}$ (resp. $\mathcal{S}'$) 表示 $Y$ (resp. $Y'$) 上的拟凝聚分次代数，即诸 $f_*(\mathcal{L}^{\otimes n})$ (resp. 诸 $f'_*(\mathcal{L}'^{\otimes n})$) 对 $n \geqslant 0$ 的直和. 注意对每个 $n \geqslant 0$，典范同态 $f'^{*}(\mathcal{S}'_n) \to \mathcal{L}'^{\otimes n}$ 的余核是 $f^{*}(\mathcal{S}_n) \to \mathcal{L}^{\otimes n}$ 的余核经 $X' \to X$ 的原像，故其支集 $Z'_n$ 是支集 $Z_n$ 的原像. 若 $\mathcal{L}'$ 丰富，诸 $Z'_n$ 的交为空，故由于 $X' \to X$ 满射，诸 $Z_n$ 的交为空，i.e. 有典范态射
\[
j \colon X \longrightarrow \mathrm{Proj}(\mathcal{S})
\]
(EGA II 3). 而且，类似态射
\[
j' \colon X' \longrightarrow \mathrm{Proj}(\mathcal{S}')
\]
正是由前者经基变换 $Y' \to Y$ 得到的 (loc. cit.). 这样，说 $\mathcal{L}'$ 相对 $f'$ 丰富意味着 $j'$ 是浸入，而且必然拟紧，因为 $f'$ 拟紧. 故由 \ref{cor:VIII.5.5} $j$ 是浸入，i.e. $\mathcal{L}$ 相对 $f$ 丰富. — 对"非常丰富"的情形完全类似地进行，上面只需取 $n = 1$，并把 $\mathrm{Proj}(\mathcal{S})$ 的考虑换成由 $\mathcal{S}_1$ 对应的射影纤维丛 $P(\mathcal{S}_1)$ 的考虑.

回忆 (EGA II 5.1.1) 拟紧态射 $f$ 称为拟仿射\index[term]{nifangshetaishe@拟仿射态射}，若对 $Y$ 中每个仿射开集 $U$，$f^{-1}(U)$ 是同构于仿射概形的开子概形的预概形. 证明 (loc. cit.) 这等价于说 $\mathcal{O}_X$ 相对 $f$ 丰富 (或也说：非常丰富). 故 \ref{cor:VIII.5.8} 蕴含：

\begin{corollary}\label{cor:VIII.5.9}
在 \ref{cor:VIII.5.4} 的条件下，且假设 $f$ 拟紧，为使 $f$ 拟仿射，必要且充分条件是 $f'$ 如此.
\end{corollary}

\begin{remark}\label{rem:VIII.5.10}
Hironaka 的非射影簇例子表明，可以有非奇异代数簇 (其中 $Y$ 射影) 的紧合态射 $f \colon X \to Y$，使 $Y$ 是两个开集 $Y_i$ 的并，$X_i = X \times_Y Y_i$ 在 $Y_i$ 上射影，但 $f$ 不射影. 故令 $Y' = Y_1 \amalg Y_2$，$Y'$ 在 $Y$ 上忠实平坦且拟紧 (甚至拟有限)，$f' \colon X' \to Y'$ 射影，但 $f$ 不射影. 因此注意，为应用 \ref{cor:VIII.5.8}，从 $f'$ 射影推出 $f$ 的同样结论，需要 $X'$ 上已经有相对 $f'$ 丰富的可逆模 $\mathcal{L}'$，配备相对 $X' \to X$ 的下降数据 (这使得可以把 $\mathcal{L}'$ 看作 $X$ 上可逆模 $\mathcal{L}$ 的逆像，后者将由 \ref{cor:VIII.5.8} 相对 $f$ 丰富). 当 $g \colon S' \to S$ 有限且局部自由时，参见 7.7.
\end{remark}

\section{应用于有限与拟有限态射}

我们要证明下列两个定理：

\begin{theorem}\label{thm:VIII.6.1}
设 $f \colon X \to Y$ 是纤维有限紧合态射\footnote{参见 EGA IV 18.12，关于对未必局部诺特预概形的推广.}，其中 $Y$ 局部诺特. 则 $f$ 有限.
\end{theorem}

\begin{theorem}\label{thm:VIII.6.2}
设 $f \colon X \to Y$ 是拟有限分离态射，其中 $Y$ 局部诺特. 则 $f$ 拟仿射，更不用说拟射影.
\end{theorem}

\begin{remark}\label{rem:VIII.6.3}
\cref{thm:VIII.6.1} 众所周知，在代数簇情形归功于 Chevalley；也见 [EGA III 4] 中用"\emph{全纯函数}定理"的简单证明. 这里给出的证明不用后者，而用下降理论；我们把它作为赠品给读者，因为得到它与 \ref{thm:VIII.6.2} 的证明"不费事". 也回忆 ([EGA III 4] 或 \cite{exp08-1}) Zariski "主定理"的全局形式，由"全纯函数定理"推出，断言：若 $f \colon X \to Y$ 拟有限且拟射影，$Y$ 诺特，则 $X$ 与有限 $Y$-预概形 $Z$ 的开子预概形 $Y$-同构. "主定理"与 \ref{thm:VIII.6.2} 的合取于是表述为：
\end{remark}

\begin{corollary}\label{cor:VIII.6.4}
设 $f \colon X \to Y$ 是拟有限分离态射，$Y$ 诺特. 则 $X$ 与有限 $Y$-预概形 $Z$ 的开子预概形 $Y$-同构.
\end{corollary}

\ref{thm:VIII.6.2} 对下降理论的另一个有趣推论将随 \ref{cor:VIII.7.9} 给出.

\begin{proof}
[6.1 与 6.2 的证明] 我们承认下列事实，其证明容易\footnote{参见 EGA IV 8.}：
\end{proof}

\begin{lemma}\label{lem:VIII.6.5}
设 $X$ 是局部诺特 $Y$ 上有限型的预概形，$y \in Y$. 为使存在 $y$ 的开邻域 $U$ 使 $X|_U$ 在 $U$ 上有限 (resp. 拟仿射，resp. ...)，必要且充分条件是 $X \times_Y \mathrm{Spec}(\mathcal{O}_y)$ 在 $\mathrm{Spec}(\mathcal{O}_y)$ 上有限 (resp. 拟仿射，resp. ...).
\end{lemma}

另一方面，由于 $f \colon X \to Y$ 有限 (resp. 拟仿射) 的性质在 $Y$ 上局部，为证 \ref{thm:VIII.6.1} 和 \ref{thm:VIII.6.2} 归结到 $Y$ 是局部环的谱，从而尤其有限维的情形. (N.B. 预概形 $Y$ 的维数定义为它的局部环的 Krull 维数的上确界). 对
\[
n = \dim(Y)
\]
作递归，$n < 0$ 时断言平凡. 故可假设 $n \geqslant 0$，且对维数 $n' < n$ 已证. 又可假设 $Y$ 是维数 $n$ 的局部诺特环 $A$ 的谱. 注意 \ref{thm:VIII.6.1} 和 \ref{thm:VIII.6.2} 中的假设随基变换保持 (在开头的归结中已用到)，它们在基变换 $\mathrm{Spec}(\widehat{A}) \to \mathrm{Spec}(A)$ 后仍然成立. 由于后者忠实平坦且拟紧，断言 \ref{cor:VIII.5.7} 和 \ref{cor:VIII.5.9} 把我们归结到 $A$ 还完备的情形. 利用事实：每个在 $A$ 上拟有限的诺特局部环 $B$ 在 $A$ 上有限，以及 $X$ 在 $Y$ 上分离且 $y$ 的纤维由孤立点组成，得到分解
\[
X = X' \amalg X''
\]
其中 $X'$ 在 $Y$ 上有限，$X''$ 在 $y$ 处的纤维为空. 若 $X$ 在 $Y$ 上紧合，则 $X''$ 亦然，故它在 $Y$ 中的像闭，且由于它不含 $y$，它为空，故 $X'' = \varnothing$ 从而 $X = X'$，这表明 $X$ 在 $Y$ 上有限并证明 \ref{thm:VIII.6.1} (N.B. 这里递归假设无用). 若 $X$ 在 $Y$ 上拟有限，$X''$ 也如此，而 $X''$ 实际位于 $Y$ 的开集 $Y - \{y\}$ 上，其维数 $< n$. 由递归假设，$X''$ 在 $Y - \{y\}$ 上拟仿射，故在 $Y$ 上也拟仿射，$X'$ 显然亦然，故它们的和 $X$ 亦然，这证明 6.2.

\begin{remark}\label{rem:VIII.6.6}
\cref{thm:VIII.6.1,thm:VIII.6.2} 在不再假设 $Y$ 局部诺特时仍然成立，只要明确假设 $f$ 有限表现 (参见 \ref{rem:VIII.3.6}). 事实上，仍可假设 $Y$ 仿射，此时容易验证情形 $f \colon X \to Y$ 由基变换 $Y \to Y_0$ 自满足与 $f$ 同样假设、$Y_0$ 诺特的情形 $f_0 \colon X_0 \to Y_0$ 得到. 故由结果 \ref{thm:VIII.6.1} resp. \ref{thm:VIII.6.2}，$f_0$ 有限 resp. 拟仿射，故 $f$ 亦然. 这类推理常在应用中有用于摆脱诺特假设 (它在应用中总是碍事).
\end{remark}

\section{下降数据的有效性判据}

如常考虑预概形态射
\[
g \colon S' \longrightarrow S
\]
和 $S'$-预概形 $X'$. 按一般事实 VII，9\footnote{不存在；参见 Séminaire Bourbaki 的 212 号报告 (1962).}，$X'$ 上相对 $g$ 的下降数据的给出，等价于等价对 \cite{exp08-3}
\[
q_1, q_2 \colon X'' \rightrightarrows X'
\]
的给出，使结构态射 $X' \to S'$ 与该对及由 $g$ 定义的等价对
\[
p_1, p_2 \colon S'' = S' \times_S S' \rightrightarrows S'
\]
相容，且使对应图式
\[
\begin{tikzcd}[row sep=large, column sep=large]
X' \arrow[d] & X'' \arrow[l] \arrow[d] \\
S' & S'' \arrow[l]
\end{tikzcd}
\]
中提取的两个方形 (或两个之一，由对称性两者等价) (用 $p_1$，$q_1$ 或 $p_2$，$q_2$) 笛卡儿. 该下降数据提出的下降问题的解，i.e. $S$ 上配备与下降数据相容的同构 $X \times_S S' \xleftarrow{\sim} X'$ 的对象 $X$，等价于满足 $h q_1 = h q_2$ 的笛卡儿方形
\[
\begin{tikzcd}[row sep=large, column sep=large]
X \arrow[d] & X' \arrow[l,"h"] \arrow[d] \\
S & S' \arrow[l,"g"]
\end{tikzcd}
\]
的给出.

由于忠实平坦拟紧态射的集随基变换稳定，且忠实平坦拟紧态射由 \ref{cor:VIII.5.3} 是有效满态射，由一般理论 \cite{exp08-D} 得出：

\begin{proposition}\label{prop:VIII.7.1}
假设 $g \colon S' \to S$ 忠实平坦且拟紧. 为使 $X'$ 上相对 $g$ 的下降数据有效，必要且充分条件是它定义的等价关系 $R = (q_1, q_2)$ 有效 (i.e. 商 $X'/R$ 存在且 $X''$ 成为 $X'$ 在 $X'/R$ 上的纤维积方形)，且典范态射 $X' \to X'/R$ 忠实平坦且拟紧.
\end{proposition}

这样，下降数据的有效性问题是一个等价图的有效性问题的特殊情形，本节给出的各种有效性判则可由此得到. 尽管如此，在下降的语境中有\cref{thm:VIII.2.1}，它蕴含若 $X'$ 在 $S'$ 上仿射，则 $X'$ 上相对 $g$ 的每个下降数据有效，这个断言对由一般平坦等价图取商没有类似物. 我们这里给出的所有有效性判则也可看作由前者推出.

设 $U'$ 是 $X'$ 的子预概形 (或更一般地 $X'$ 在范畴 $\mathbf{Sch}$ 中的子对象)；称 $U'$ 被 $X'$ 上的下降数据稳定，若可以在 $U'$ 上找到相对 $g$ 的下降数据，使浸入 $U' \to X'$ 与下降数据相容. 这也意味着 $U'$ 在 $X''$ 中经 $q_1$ 和 $q_2$ 的逆像相同 (或如常所说的，$U'$ 被等价关系 $R$ 稳定)，当然 $U'$ 上的所求下降数据于是唯一，称为由 $X'$ 的下降数据诱导的下降数据. 这样：

\begin{proposition}\label{prop:VIII.7.2}
设 $(X'_i)$ 是 $X'$ 的被下降数据稳定的开集 $X'_i$ 的覆盖. 为使 $X'$ 上的下降数据有效，必要且充分条件是诸 $X'_i$ 中诱导的下降数据有效.
\end{proposition}

这是例如 \ref{prop:VIII.7.1} 的容易推论，证明细节留给读者.

\begin{corollary}\label{cor:VIII.7.3}
设 $(S_i)$ 是 $S$ 的开覆盖，对每个 $i$ 设 $S'_i$ 和 $X'_i$ 由 $S'$ 和 $X'$ 经基变换 $S_i \to S$ 得到. 为使 $X'$ 上的下降数据有效，必要且充分条件是：对每个 $i$，$X'_i$ 上相对 $g_i \colon S'_i \to S_i$ 的下降数据有效.
\end{corollary}

这个判则实际上总把我们归结到 $S$ 仿射的情形. 在 $S'$ 也仿射的情形 (应用中最常见)，有：

\begin{corollary}\label{cor:VIII.7.4}
假设 $S$ 和 $S'$ 仿射. 为使 $X'$ 上的下降数据有效，必要且充分条件是 $X'$ 是被下降数据稳定的仿射开集 $X'_i$ 的并.
\end{corollary}

充分性来自 \ref{prop:VIII.7.2} 及事实：若 $X'_i$ 仿射，它在 $S'$ 上仿射，可应用 2.1. 对必要性，注意若 $X'$ 来自 $X$，且 $X$ 被仿射开集 $X_i$ 覆盖，则诸 $X'_i = X_i \times_S S'$ 是被下降数据稳定且覆盖 $X'$ 的仿射开集.

\begin{corollary}\label{cor:VIII.7.5}
设 $g \colon S' \to S$ 是忠实平坦、拟紧且根式的态射. 则 $g$ 是有效下降态射，i.e. 对 $S'$ 上的每个 $X'$，$X'$ 上相对 $g \colon S' \to S$ 的每个下降数据有效.
\end{corollary}

事实上，由 \ref{cor:VIII.7.3} 可假设 $S$ 仿射，故由于 $S'$ 在 $S$ 上根式故分离，$S'$ 分离. 而且对每个 $x' \in X'$，由等价关系 $R$ 定义的集合等价关系的纤维 $R(x') = q_2(q_1^{-1}(x'))$ 约化为一点，因为 $g$ 根式时，由它经基变换 $S' \to S$ 得到的 $p_1$，$p_2$ 也根式，故由前者经基变换 $X' \to S''$ 得到的 $q_1$，$q_2$ 也根式. 故 $X'$ 的每个开集被下降数据稳定. 用仿射开集 $X'_i$ 覆盖 $X'$，由于 $S'$ 分离它们在 $S$ 上仿射，故由 \ref{thm:VIII.2.1} 诱导的下降数据有效. 再由 \ref{prop:VIII.7.2} 得出结论.

注意 \ref{cor:VIII.7.5} 给出预概形范畴中已知的有效下降态射的唯一情形，而且这很可能确实是唯一情形，即使限于诺特概形或域上有限型概形. 当假设 $S$ 局部诺特且 $S'$ 在 $S$ 上有限型时，断言 \ref{cor:VIII.7.5} 也是下列 (推广 Weil 的 Galois 下降和 Cartier 的不可分下降) 的特殊情形：

\begin{corollary}\label{cor:VIII.7.6}
设 $g \colon S' \to S$ 是有限局部自由 (i.e. 由 $S$ 上为有限型局部自由模的代数定义) 且满射 (故 $g$ 忠实平坦且拟紧，从而是下降态射) 的态射. 设 $X'$ 是配备下降数据的 $S'$-预概形. 为使该数据有效，必要且充分条件是：对每个 $x' \in X'$，纤维 $R(x') = q_2(q_1^{-1}(x'))$ 包含在某个仿射开集中 (若 $X'$ 在 $S'$ 上拟射影，该条件自动满足).
\end{corollary}

括号中的评注来自事实：若 $s$ 是 $x'$ 下方的 $S$ 的点，则 $R(x')$ 有限且包含在纤维 $X'_s$ 中，另一方面由于 $X'$ 在 $S'$ 上拟射影且 $S'$ 在 $S$ 上有限，$X'$ 在 $S$ 上拟射影，这意味着 $X'$ 在 $S$ 上的纤维包含在仿射开集中.

由于仿射概形的任何有限部分有仿射邻域基本系，看到限制到 $S$ 的某个仿射开集上方不损失假设，由 \ref{cor:VIII.7.3} 这归结到 $S$ 仿射的情形. 由 \ref{cor:VIII.7.4}，归结到证明 $x'$ 包含在被下降数据稳定的仿射开集中. 设 $U$ 是包含 $R(x')$ 的仿射开集，则饱和集
\[
R(X' - U) = q_2(q_1^{-1}(X' - U))
\]
与 $R(x')$ 不相交，另一方面由于 $q_2$ 有限 (因为 $g$ 故 $p_2$ 有限) 故闭，第二个成员是 $X'$ 的闭部分. 设 $U'$ 是它在 $X'$ 中的余集，它是饱和开集，且
\[
R(x') \subset U' \subset U
\]
其中 $U$ 仿射，但 $U'$ 事先未必仿射. 由于仿射概形 $U$ 中的有限部分 $R(x')$ 有 $U_f$ 形的仿射邻域基本系，把 $f$ 换成它在 $U'$ 上的限制，看到存在 $\mathcal{O}_U$ 的截面 $f$ 使：
\[
R(x') \subset U'_f, \quad U'_f \text{ 仿射}.
\]
设 $U'' = q_1^{-1}(U') = q_2^{-1}(U')$，仍用 $q_1$，$q_2$ 记诱导态射 $U'' \to U'$，考虑
\[
f' = N_{q_2}(q_1^{*}(f)),
\]
其中 $N_{q_2}$ 表示相对有限局部自由态射 $q_2 \colon U'' \to U'$ 的范数. 范数构造与基变换的相容性容易蕴含 $f'$ 是不变截面：
\[
q_1^{*}(f') = q_2^{*}(f')
\]
这蕴含 $U'_{f'}$ 是 $U'$ 的饱和开集. 而且更精确地，用 $Z(f')$ 记截面 $f'$ 的零点集，由范数的性质得到：
\[
Z(f') = q_2(Z(q_1^{*}(f))) = q_2(q_1^{*}(Z(f))) = R(U' - U'_f),
\]
这蕴含 $U'_{f'} = U' - Z(f')$ 饱和，包含 $R(x')$，且包含在 $U'_f$ 中. 由于后者仿射，得出 $U'_{f'}$ 也仿射 (因为它等于 $(U'_f)_{f''}$，其中 $f'' = f'|_{U'_{f'}}$). 故它是包含 $R(x')$ 从而 $x'$ 的仿射饱和开集，这完成证明.

注意这个推理适用于预概形 $X'$ 中有限局部自由的任何等价关系 (或甚至只是\emph{预等价}\index[term]{yudengjia@预等价}，参见 \cite{exp08-3})，而且 \ref{cor:VIII.7.6} 也是有限局部自由预等价的类似结果 (loc. cit.) 的特殊情形. 对下面的 \ref{cor:VIII.7.7} 同样评注. 也可以在得到包含 $x'$ 的仿射 (拟仿射) 饱和开集 $U'$ 的存在后，借助 \ref{cor:VIII.7.9} 和 \ref{prop:VIII.7.2}，这避免范数的使用. 还注意在 \ref{cor:VIII.7.6} 的条件下，若 $X'$ 上的下降数据有效，$X'$ 来自 $S$ 上的 $X$，则态射 $X' \to X$ 有限、局部自由且满射，因为它由 $g$ 经基变换 $X \to S$ 得到. 由此得出 (EGA II 6.6.4) 若 $X'$ 在 $S'$ 上故在 $S$ 上拟射影，则 $X$ 在 $S$ 上拟射影 (取 $X$ 上相对丰富可逆层为 $X'$ 上相对丰富 (于 $S$，或于 $S'$，两者相同) 可逆层的范数). 由此得到：

\begin{corollary}\label{cor:VIII.7.7}
有限局部自由且满射的态射 $g \colon S' \to S$ 对互相拟射影的预概形的纤维范畴是有效下降态射，i.e. 对每个在 $S'$ 上拟射影的 $X'$，$X'$ 上相对 $g$ 的每个下降数据有效，且下降得到的 $S$-预概形 $X$ 在 $S$ 上拟射影.
\end{corollary}

\begin{proposition}\label{prop:VIII.7.8}
设 $g \colon S' \to S$ 是忠实平坦拟紧态射. 则 $g$ 对预概形 $T$ 上配备相对 $T$ 丰富可逆层的拟紧预概形 $Z$ 的纤维范畴是有效下降态射. 特别，对 $S'$ 上的每个预概形 $X'$，配备相对 $g \colon S' \to S$ 的下降数据，以及对 $X'$ 上相对 $S'$ 丰富的每个可逆层 $\mathcal{L}'$，配备相对 $X'$ 上给定数据的下降数据，(i.e. 配备满足通常传递条件的 $q_1^{*}(\mathcal{L}')$ 与 $q_2^{*}(\mathcal{L}')$ 的同构)，$X'$ 上的下降数据有效，且由下降自 $\mathcal{L}'$ 得到的下降预概形 $X$ 上的可逆层 $\mathcal{L}$ 相对 $S$ 丰富.
\end{proposition}

证明与 \ref{cor:VIII.5.8} 完全类似，注意由 $\mathcal{L}'$ 定义的 $S'$ 上拟凝聚分次代数 $\mathcal{S}'$ 上有下降数据，使得可以借助 \ref{thm:VIII.1.1} 构造 $S$ 上拟凝聚分次代数 $\mathcal{S}$，从而得到 $S$ 上的 $P = \mathrm{Proj}(\mathcal{S})$，使 $P' = \mathrm{Proj}(\mathcal{S}')$ 连同其下降数据等同于 $P \times_S S'$. 由于由假设 $X'$ 等同于 $P'$ 的开集 (必然被 $P'$ 上的下降数据稳定)，$X'$ 上的下降数据也有效，得到的下降预概形是 $P$ 中的开集. 细节留给读者. — 特别取 $\mathcal{L}' = \mathcal{O}_{X'}$，得到：

\begin{corollary}\label{cor:VIII.7.9}
设 $g \colon S' \to S$ 是忠实平坦拟紧态射，设 $X'$ 是 $S'$ 上方的拟仿射预概形，则 $X'$ 上相对 $g$ 的每个下降数据有效，且下降得到的预概形 $X$ 在 $S$ 上拟仿射.
\end{corollary}

由 \ref{thm:VIII.6.2}，这个结果特别适用于 $S'$ 局部诺特且 $X'$ 在 $S'$ 上拟有限分离的情形，更一般地适用于 $S'$ 任意且 $X'$ 在 $S'$ 上有限表现、拟有限分离的情形 (参见 \ref{rem:VIII.6.6}).

\begin{remark}\label{rem:VIII.7.10}
本节给出的结果穷尽了目前已知的有效性判则，很可能甚至穷尽了所有有用的判则\footnote{这种看法部分被证明有误，参见 J.P. Murre，Séminaire Bourbaki 294 (附录)，1965年5月，以及一些特殊结果 (特别是 Néron 和 Raynaud 的). 关于群概形的下降，参见 M. Raynaud 的博士论文 (1968).}. 支持这一断言的注意下列反例：
\begin{enumerate}[i)]
\item 若 $S$ 是域的谱，$S'$ 是二次 伽罗瓦扩张的谱，可以找到 $S'$ 上紧合且光滑、维数 3 的 $X'$，配备不有效的下降数据 (Serre).
\item 可以找到 $S$ 为维数 3 的正则局部环的谱 (若愿意，可以取给定特征域上代数概形的局部环)，$T$ 是 $S$ 的群 $Z/2Z$ 的主覆盖，使若 $t$ 表示 $T$ 的位于 $S$ 的闭点 $s$ 上方的点之一，且 $S' = T - s$，则能找到 $S'$ 上射影且正则的 $X'$，配备相对 $g \colon S' \to S$ 的下降数据，该下降数据不有效.
\end{enumerate}
这些构造使用 Hironaka 的非射影簇例子. 对 (i)，只需用事实：可以在 $k$ 上方找到维数 3 的紧合光滑概形 $X_0$，$G = Z/2Z$ 无惯性地作用其上，其中有 $S$ 上两个有理点 $a$，$b$，在 $G$ 下相伴，不包含在仿射开集中. 令 $X' = X_0 \times_k k'$，借助 $G$ 在两个因子上的作用使 $G$ 作用在 $X'$ 上，这给出 $X'$ 上相对 $g \colon \mathrm{Spec}(k') \to \mathrm{Spec}(k)$ 的下降数据. 在 $a$ resp. $b$ 上方恰好有一个点 $a'$ resp. $b'$ (带二次剩余扩张)，且 $a'$ 和 $b'$ 在 $G$ 下相伴，因为 $X' \to X_0$ 与 $G$ 的作用相容. 则 $a'$ 和 $b'$ 不能包含在仿射开集 $U'$ 中，因为此时 $U = X_0 - \mathrm{Im}(X' - U')$ 将是包含 $(a, b)$ 且逆像在 $X'$ 中位于 $U'$ 内的开集，故拟仿射，从而 $U$ 拟仿射，于是 $(a, b)$ 在 $U$ 中有仿射邻域.

对 (ii)，用事实：在 Hironaka 的例子中，$X_0$ 作为 $k$-射影概形 $Y$ ($k$ 上光滑，态射 $f \colon X_0 \to Y$ 双有理，但无关紧要) 上方的紧合预概形得到，群 $G$ 也作用在 $Y$ 上，与它在 $X_0$ 上的作用相容，最后令 $S' = Y - f(b)$，$X' = X_0|_{S'}$，$X'$ 在 $S'$ 上射影. 则 $X_0$ 配备相对典范态射 $Y \to S = Y/G$ 的自然下降数据，借助 $G$ 在 $X_0$ 上与它在 $Y$ 上的作用相容的作用. 该下降数据不有效，因为 $(a, b)$ 不包含在仿射开集中. $X'$ 上相对 $g \colon S' \to S$ 的诱导下降数据于是不有效，容易验证.
\end{remark}

\section{文献}\printbibliography[heading=none]


